% flow_hamilton_variational.tex — extended Hamilton principle: energy/variational structure
% One functional Pi -> stationarity delta Pi = 0 -> mechanics (left) & material-point (right).
% Usage: \resizebox{\linewidth}{!}{% flow_hamilton_variational.tex — extended Hamilton principle: energy/variational structure
% One functional Pi -> stationarity delta Pi = 0 -> mechanics (left) & material-point (right).
% Usage: \resizebox{\linewidth}{!}{% flow_hamilton_variational.tex — extended Hamilton principle: energy/variational structure
% One functional Pi -> stationarity delta Pi = 0 -> mechanics (left) & material-point (right).
% Usage: \resizebox{\linewidth}{!}{% flow_hamilton_variational.tex — extended Hamilton principle: energy/variational structure
% One functional Pi -> stationarity delta Pi = 0 -> mechanics (left) & material-point (right).
% Usage: \resizebox{\linewidth}{!}{\input{flow_hamilton_variational.tex}}
\begin{tikzpicture}[
  every node/.style = {font=\small},
  N/.style  = {draw, rounded corners=3pt, align=center, thick,
               text width=5.0cm, minimum height=0.85cm, inner sep=4pt},
  Ntop/.style= {N, fill=purple!12, text width=8.6cm},
  Nen/.style = {N, fill=cyan!12},
  Nop/.style = {draw, ellipse, align=center, thick, fill=yellow!22,
                inner sep=3pt, font=\small, text width=4.6cm},
  Nmech/.style={N, fill=orange!15, text width=5.4cm},
  Nbio/.style ={N, fill=teal!14, text width=5.4cm},
  ar/.style = {-Stealth, thick},
  lbl/.style= {font=\scriptsize\itshape, fill=white, inner sep=1.5pt, align=center},
]

\node[Ntop] (pi) {%
  \textbf{Total potential functional} $\Pi[\mathbf u,\bm\phi,\alpha]$\\[1pt]
  {\scriptsize the system's total (free) energy — one scalar}};

\node[Nen, below left=1.2cm and 0.2cm of pi] (el) {%
  \textbf{elastic strain energy}\\ $\Pi_{\mathrm{el}}(\mathbf F_e)$};
\node[Nen, below=1.2cm of pi] (ch) {%
  \textbf{chemical / growth potential}\\ $\Pi_{\mathrm{ch}}(\bm\phi,\alpha)$};
\node[Nen, below right=1.2cm and 0.2cm of pi] (di) {%
  \textbf{dissipation}\\ $\mathcal{D}$ (interface damage, viscosity)};

\node[Nop, below=1.15cm of ch] (var) {%
  first variation / stationarity\\ $\delta\Pi=0$\;\;(one operation)};

\node[Nmech, below left=1.15cm and -0.2cm of var] (mech) {%
  \textbf{Mechanical balance}\\[1pt]
  $\nabla\!\cdot\bm\sigma=\mathbf 0$\;
  (Cauchy stress, $\mathbf F=\mathbf F_e\mathbf F_g$)};
\node[Nbio, below right=1.15cm and -0.2cm of var] (bio) {%
  \textbf{Material-point equations}\\[1pt]
  growth $\dot\alpha=k_\alpha\phi$ \& cross-diffusion evolution of $\bm\phi$};

\draw[ar] (pi.south) -- ++(0,-0.3) coordinate (s1);
\draw[ar] (s1) -| (el.north);
\draw[ar] (s1) -- (ch.north);
\draw[ar] (s1) -| (di.north);
\draw[ar] (el.south) |- (var.west);
\draw[ar] (ch.south) -- (var.north);
\draw[ar] (di.south) |- (var.east);
\draw[ar] (var.south) -- ++(0,-0.3) coordinate (s2);
\draw[ar] (s2) -| (mech.north) node[lbl,pos=0.25,above]{$\delta_{\mathbf u}\Pi$};
\draw[ar] (s2) -| (bio.north)  node[lbl,pos=0.25,above]{$\delta_{\bm\phi,\alpha}\Pi$};

\node[font=\scriptsize\itshape, text=gray!55, align=center, text width=8cm]
      at ($(mech.south)!0.5!(bio.south) + (0,-0.45)$)
  {mechanics and biology fall out of the \emph{same} energy — simultaneously};

\end{tikzpicture}
}
\begin{tikzpicture}[
  every node/.style = {font=\small},
  N/.style  = {draw, rounded corners=3pt, align=center, thick,
               text width=5.0cm, minimum height=0.85cm, inner sep=4pt},
  Ntop/.style= {N, fill=purple!12, text width=8.6cm},
  Nen/.style = {N, fill=cyan!12},
  Nop/.style = {draw, ellipse, align=center, thick, fill=yellow!22,
                inner sep=3pt, font=\small, text width=4.6cm},
  Nmech/.style={N, fill=orange!15, text width=5.4cm},
  Nbio/.style ={N, fill=teal!14, text width=5.4cm},
  ar/.style = {-Stealth, thick},
  lbl/.style= {font=\scriptsize\itshape, fill=white, inner sep=1.5pt, align=center},
]

\node[Ntop] (pi) {%
  \textbf{Total potential functional} $\Pi[\mathbf u,\bm\phi,\alpha]$\\[1pt]
  {\scriptsize the system's total (free) energy — one scalar}};

\node[Nen, below left=1.2cm and 0.2cm of pi] (el) {%
  \textbf{elastic strain energy}\\ $\Pi_{\mathrm{el}}(\mathbf F_e)$};
\node[Nen, below=1.2cm of pi] (ch) {%
  \textbf{chemical / growth potential}\\ $\Pi_{\mathrm{ch}}(\bm\phi,\alpha)$};
\node[Nen, below right=1.2cm and 0.2cm of pi] (di) {%
  \textbf{dissipation}\\ $\mathcal{D}$ (interface damage, viscosity)};

\node[Nop, below=1.15cm of ch] (var) {%
  first variation / stationarity\\ $\delta\Pi=0$\;\;(one operation)};

\node[Nmech, below left=1.15cm and -0.2cm of var] (mech) {%
  \textbf{Mechanical balance}\\[1pt]
  $\nabla\!\cdot\bm\sigma=\mathbf 0$\;
  (Cauchy stress, $\mathbf F=\mathbf F_e\mathbf F_g$)};
\node[Nbio, below right=1.15cm and -0.2cm of var] (bio) {%
  \textbf{Material-point equations}\\[1pt]
  growth $\dot\alpha=k_\alpha\phi$ \& cross-diffusion evolution of $\bm\phi$};

\draw[ar] (pi.south) -- ++(0,-0.3) coordinate (s1);
\draw[ar] (s1) -| (el.north);
\draw[ar] (s1) -- (ch.north);
\draw[ar] (s1) -| (di.north);
\draw[ar] (el.south) |- (var.west);
\draw[ar] (ch.south) -- (var.north);
\draw[ar] (di.south) |- (var.east);
\draw[ar] (var.south) -- ++(0,-0.3) coordinate (s2);
\draw[ar] (s2) -| (mech.north) node[lbl,pos=0.25,above]{$\delta_{\mathbf u}\Pi$};
\draw[ar] (s2) -| (bio.north)  node[lbl,pos=0.25,above]{$\delta_{\bm\phi,\alpha}\Pi$};

\node[font=\scriptsize\itshape, text=gray!55, align=center, text width=8cm]
      at ($(mech.south)!0.5!(bio.south) + (0,-0.45)$)
  {mechanics and biology fall out of the \emph{same} energy — simultaneously};

\end{tikzpicture}
}
\begin{tikzpicture}[
  every node/.style = {font=\small},
  N/.style  = {draw, rounded corners=3pt, align=center, thick,
               text width=5.0cm, minimum height=0.85cm, inner sep=4pt},
  Ntop/.style= {N, fill=purple!12, text width=8.6cm},
  Nen/.style = {N, fill=cyan!12},
  Nop/.style = {draw, ellipse, align=center, thick, fill=yellow!22,
                inner sep=3pt, font=\small, text width=4.6cm},
  Nmech/.style={N, fill=orange!15, text width=5.4cm},
  Nbio/.style ={N, fill=teal!14, text width=5.4cm},
  ar/.style = {-Stealth, thick},
  lbl/.style= {font=\scriptsize\itshape, fill=white, inner sep=1.5pt, align=center},
]

\node[Ntop] (pi) {%
  \textbf{Total potential functional} $\Pi[\mathbf u,\bm\phi,\alpha]$\\[1pt]
  {\scriptsize the system's total (free) energy — one scalar}};

\node[Nen, below left=1.2cm and 0.2cm of pi] (el) {%
  \textbf{elastic strain energy}\\ $\Pi_{\mathrm{el}}(\mathbf F_e)$};
\node[Nen, below=1.2cm of pi] (ch) {%
  \textbf{chemical / growth potential}\\ $\Pi_{\mathrm{ch}}(\bm\phi,\alpha)$};
\node[Nen, below right=1.2cm and 0.2cm of pi] (di) {%
  \textbf{dissipation}\\ $\mathcal{D}$ (interface damage, viscosity)};

\node[Nop, below=1.15cm of ch] (var) {%
  first variation / stationarity\\ $\delta\Pi=0$\;\;(one operation)};

\node[Nmech, below left=1.15cm and -0.2cm of var] (mech) {%
  \textbf{Mechanical balance}\\[1pt]
  $\nabla\!\cdot\bm\sigma=\mathbf 0$\;
  (Cauchy stress, $\mathbf F=\mathbf F_e\mathbf F_g$)};
\node[Nbio, below right=1.15cm and -0.2cm of var] (bio) {%
  \textbf{Material-point equations}\\[1pt]
  growth $\dot\alpha=k_\alpha\phi$ \& cross-diffusion evolution of $\bm\phi$};

\draw[ar] (pi.south) -- ++(0,-0.3) coordinate (s1);
\draw[ar] (s1) -| (el.north);
\draw[ar] (s1) -- (ch.north);
\draw[ar] (s1) -| (di.north);
\draw[ar] (el.south) |- (var.west);
\draw[ar] (ch.south) -- (var.north);
\draw[ar] (di.south) |- (var.east);
\draw[ar] (var.south) -- ++(0,-0.3) coordinate (s2);
\draw[ar] (s2) -| (mech.north) node[lbl,pos=0.25,above]{$\delta_{\mathbf u}\Pi$};
\draw[ar] (s2) -| (bio.north)  node[lbl,pos=0.25,above]{$\delta_{\bm\phi,\alpha}\Pi$};

\node[font=\scriptsize\itshape, text=gray!55, align=center, text width=8cm]
      at ($(mech.south)!0.5!(bio.south) + (0,-0.45)$)
  {mechanics and biology fall out of the \emph{same} energy — simultaneously};

\end{tikzpicture}
}
\begin{tikzpicture}[
  every node/.style = {font=\small},
  N/.style  = {draw, rounded corners=3pt, align=center, thick,
               text width=5.0cm, minimum height=0.85cm, inner sep=4pt},
  Ntop/.style= {N, fill=purple!12, text width=8.6cm},
  Nen/.style = {N, fill=cyan!12},
  Nop/.style = {draw, ellipse, align=center, thick, fill=yellow!22,
                inner sep=3pt, font=\small, text width=4.6cm},
  Nmech/.style={N, fill=orange!15, text width=5.4cm},
  Nbio/.style ={N, fill=teal!14, text width=5.4cm},
  ar/.style = {-Stealth, thick},
  lbl/.style= {font=\scriptsize\itshape, fill=white, inner sep=1.5pt, align=center},
]

\node[Ntop] (pi) {%
  \textbf{Total potential functional} $\Pi[\mathbf u,\bm\phi,\alpha]$\\[1pt]
  {\scriptsize the system's total (free) energy — one scalar}};

\node[Nen, below left=1.2cm and 0.2cm of pi] (el) {%
  \textbf{elastic strain energy}\\ $\Pi_{\mathrm{el}}(\mathbf F_e)$};
\node[Nen, below=1.2cm of pi] (ch) {%
  \textbf{chemical / growth potential}\\ $\Pi_{\mathrm{ch}}(\bm\phi,\alpha)$};
\node[Nen, below right=1.2cm and 0.2cm of pi] (di) {%
  \textbf{dissipation}\\ $\mathcal{D}$ (interface damage, viscosity)};

\node[Nop, below=1.15cm of ch] (var) {%
  first variation / stationarity\\ $\delta\Pi=0$\;\;(one operation)};

\node[Nmech, below left=1.15cm and -0.2cm of var] (mech) {%
  \textbf{Mechanical balance}\\[1pt]
  $\nabla\!\cdot\bm\sigma=\mathbf 0$\;
  (Cauchy stress, $\mathbf F=\mathbf F_e\mathbf F_g$)};
\node[Nbio, below right=1.15cm and -0.2cm of var] (bio) {%
  \textbf{Material-point equations}\\[1pt]
  growth $\dot\alpha=k_\alpha\phi$ \& cross-diffusion evolution of $\bm\phi$};

\draw[ar] (pi.south) -- ++(0,-0.3) coordinate (s1);
\draw[ar] (s1) -| (el.north);
\draw[ar] (s1) -- (ch.north);
\draw[ar] (s1) -| (di.north);
\draw[ar] (el.south) |- (var.west);
\draw[ar] (ch.south) -- (var.north);
\draw[ar] (di.south) |- (var.east);
\draw[ar] (var.south) -- ++(0,-0.3) coordinate (s2);
\draw[ar] (s2) -| (mech.north) node[lbl,pos=0.25,above]{$\delta_{\mathbf u}\Pi$};
\draw[ar] (s2) -| (bio.north)  node[lbl,pos=0.25,above]{$\delta_{\bm\phi,\alpha}\Pi$};

\node[font=\scriptsize\itshape, text=gray!55, align=center, text width=8cm]
      at ($(mech.south)!0.5!(bio.south) + (0,-0.45)$)
  {mechanics and biology fall out of the \emph{same} energy — simultaneously};

\end{tikzpicture}
